% Transcription — fonds Alexandre Grothendieck, shelfmark 156-9, batch 2 (pages 21-40).
% Established from the facsimile archives/batches/156-9/batch-02.pdf, per the
% skill transcribe-grothendieck.
%
% Pass: Opus 5.5 (claude-opus-5-5), 2026-10-03 — first pass, unchecked against the pages by a human.
%
% A fast drafting hand throughout: the formulas carry, much of the connective
% prose does not, and the apparatus is correspondingly heavy.
\documentclass[11pt,a4paper]{article}
\input{../preamble/grothendieck}

\folder{156-9}
\batch{2}
\pages{21}{40}
\foldertitle{[Chapitre] IX et IX bis. [Ateliers] : notes manuscrites (05-15/07/1986).}
\dating{1986}
\watermark{Édition de démonstration}

\begin{document}

\page{21}
\note{p. 18 de l'auteur ; le calcul sur $\complement\Phi$ commence avant ce lot.}

\[
  \underbrace{(\complement\Phi)^{**}}_{\substack{\text{polygone}\\ \text{induit}\\ \text{associé à } \complement\Phi}} \;=\; \Phi^{*}
\]
\note{la légende sous l'accolade est de lecture incertaine (« polygone induit associé à $\complement\Phi$ »).}

d'où, en tenant compte de la formule fondamentale \add{plus haut}

\[
  \boxed{\operatorname{Cosupp}^{\circ}\Phi \;=\; \operatorname{supp}^{\circ}\operatorname{Omb}^{\circ}\complement\Phi \;=\; \operatorname{supp}^{\circ}\complement\Phi .}
\]
\note{sous la ligne encadrée, un second membre « $=\operatorname{cosupp}^{\circ}\ldots$ » est biffé ; le premier $\operatorname{supp}^{\circ}$ est récrit sur un mot biffé illisible.}

Soit alors, pour tout $\Phi\subset\Phi_T$,
\[
  S_\Phi\subset\mathcal{M}, \qquad S_\Phi=\operatorname{supp}^{\circ}(\Omega) =\operatorname{supp}^{\circ}(\operatorname{Omb}^{\circ}\Phi) .
\]
\note{dans le second membre, l'argument de $\operatorname{supp}^{\circ}$ est surchargé et biffé, puis récrit « $\operatorname{supp}^{\circ}(\operatorname{Omb}^{\circ}\Phi)$ » ; le $\Omega$ ci-dessus ne rend pas une lecture, il tient la place de ce premier argument, qu'on lit sans doute $\Phi$.}

On trouve une application
\[
  \Phi\longmapsto S_\Phi, \qquad \mathfrak{P}(\Phi_T)\longrightarrow \underbrace{\Sigma_{\mathrm{cons}}(\mathcal{M})}_{\substack{\text{ens. des}\\ \text{supports}\\ \text{constructibles}}}
\]
Cette application transforme \uncertain{(par définition)} \ill{} sommes (i.e. $\Sigma_\Phi$ pour $c$) en $\operatorname{Sup}$ dans $\Sigma_{\mathrm{cons}}$ \add{de $\mathcal{M}$}. La formule précédente signifie qu'elles commutent à $\complement$. Il en résulte qu'elles commutent aux intersections. \uncertain{Cela utilise} l'axiome des supports Ax\,supp\,1.

Deux \struck{parties} polygones \add{$\Phi,\Psi$} sont dits \textit{commensurables} si $\delta\Phi\cup\delta\Psi$ est un diagramme. Cela \struck{signifie} explique que leurs

\page{22}
\note{p. 19 de l'auteur.}
supports \add{$\operatorname{supp}^{\circ}$} sont \textit{« subordonnés »} à un même diagramme $\delta\Phi\cup\delta\Psi\cup\partial I=T$. [Il en \uncertain{est ainsi} pour un ensemble quelconque de polygones, \ill{} de partitions constructibles de $\mathcal{M}$ (\uncertain{cf. ex. en appendice}).] \uncertain{cqfd} — quitte à subdiviser les intervalles par polygones dans $\Phi$, $\Psi$, ce qui ne change pas les \uncertain{$\operatorname{supp}^{\circ}$}).

Toutes les opérations \add{usuelles} \ill{} qu'on peut faire, sur une famille quelconque \uncertain{d'éléments} de $\Sigma_{\mathrm{cons}}$ mutuellement commensurables, se ramènent aux opérations $\cup,\cap,\complement$ dans un $\mathfrak{P}(\Phi_T)$ convenable.

NB Si $T$ \uncertain{de longueur} $n$, ou $n$ intervalles et \uncertain{$n+1$ sommets} ($n\geq 1$),
\[
  \operatorname{card}\Phi_T=2n+1
\]
et
\[
  \underbrace{\operatorname{card}\mathfrak{P}(\Phi_T)}_{\substack{\text{nb de supports}\\ \text{subordonnés à } T}} \;=\; 2^{2n+1}.
\]

\textbf{2)} Cas d'un $I$ où on ne se donne pas $\omega_0,\omega_1$. On doit le voir comme se déduisant du cas précédent $(I,\partial I)$, en prenant un

\page{23}
\note{p. 20 de l'auteur.}
« origine virtuelle » ou « supplémentaire » p.r. à une des extrémités de $I$, ou des deux. Supposons p.\,ex. \struck{qu'il n'y ait} qu'il \ill{} \uncertain{origine} à $(I,\partial I)\smallsetminus\{\omega_0\}$ \ill{} \struck{plus} \ill{} \uncertain{partie}) — ou, plus précisément, le \uncertain{commencement}
\[
  \text{\struck{\ill{}}}\ \operatorname{Omb}^{\circ}(\{I_{\omega_0,\omega_1},\omega_1\})=\mathcal{M}\smallsetminus\omega_0
\]
(partie de $\mathcal{M}$ fermée \uncertain{par} $\omega_0$), on \uncertain{décrit} les \ill{} parties en termes de
\[
  (\overbrace{I\smallsetminus\{\omega_0\}}^{J},\ \omega_1)
\]
un ordonné $J$, \uncertain{muni} d'un plus grand élément $\omega_1$. On garde \struck{tous} les strates de $\mathcal{M}(I,\partial I)$, à l'exception \add{\ill{}} de $\omega_0$. Mais les strates de la forme $I_{\omega_0,a}$ \struck{s'identifient} ($\omega_0<a$ i.e. $a\in J$) \uncertain{maintenant} $J_{<a}$.
\marginal{On ne suppose plus ici que $a$ a un seul \uncertain{prédécesseur}, cf. \uncertain{plus bas}}

Les relations $\trianglelefteq$, $\leq$, $\mathring{\ll}$, $|0|$ \struck{et} \uncertain{dans} l'ens.
\[
  \mathcal{M}(J,\omega_1)\ (=\mathcal{M}(I,\partial I)\smallsetminus\{\omega_0\})
\]
\uncertain{peuvent alors s'expliciter} \struck{\ill{}} \add{explicit\supplied{ement}}, \uncertain{en} termes de $J,\omega_1$, \uncertain{sans} réf.\ à $(I,\partial I)$, \uncertain{en travaillant} avec les strates
\[
  a,\quad J_{<a},\quad J_{a,b}\qquad a<b,\ a,b\in J .
\]
Si p.\,ex.\ $I=\{\omega_0,\omega_1\}$, on aura $J=\{\omega_1\}$,

\page{24}
\note{p. 21 de l'auteur.}
et
\[
  \mathcal{M}(J,\omega_1)=\{\omega_1,\ J_{<\omega_1}\}
\]
magasin qui consiste en \textit{deux} éléments seulement, et \uncertain{i.e.} \textit{\uncertain{une maquette}}. Ici, on a
\[
  |J_{<\omega_1}|^{\circ}=\emptyset, \quad\text{mais bien sûr}\quad \mathring{J}_{<\omega_1}\neq\emptyset .
\]

\struck{Choix \ill{}} Traitement \uncertain{dual}

\[
  \text{pour}\qquad \mathcal{M}(I,\partial I)\smallsetminus\{\omega_1\} \overset{\mathrm{def}}{=} \mathcal{M}(J,\omega_0)
\]
\uncertain{où} $\omega_0$ est désormais un \struck{tout} \uncertain{tel que} plus petit élément de $J=I\smallsetminus\{\omega_1\}$. Cette fois, \uncertain{ce sont} les éléments
\[
  J_{>a}\qquad (a\in J)
\]
qu'il faut introduire, en plus des « \uncertain{lieux} » $a\in J$ et des « intervalles finis » $J_{a,b}$.

\uncertain{Prenons} enfin le cas des magasins \uncertain{privés} de $\partial I$ dans $\mathcal{M}$
\[
  \mathcal{M}(I,\partial I)\smallsetminus\partial I \overset{\mathrm{def}}{=} \mathcal{M}(J) \qquad J=I\smallsetminus\partial I .
\]
Cette fois, il faut introduire à la fois les $J_{<a}$ et les $J_{>a}$ ($a\in J$),
\marginal{\ill{} \uncertain{quand bien même} $J$ ! \ill{}}
\uncertain{remplaçant} les « traces » des intervalles \struck{\ill{}} $I_{\omega_0,a}$, $I_{a,\omega_1}$, et \struck{\ill{}} enfin, de plus, une \uncertain{strate} unité
\[
  [J]
\]
qui figure \uncertain{la trace} des $I_{\omega_0,\omega_1}$. L'introduction

\page{25}
\note{p. 22 de l'auteur.}
est indispensable, si on veut pouvoir exprimer tous les \uncertain{\ill{}} d'une figure complète $\Phi_T$ ($T\in\operatorname{Diag}(J)$) telles que
\[
  \{J_{<a},\ J_{>a}\}\quad\text{si } a\in J
\]
ou
\[
  J_{<a},\ J_{a,b},\ J_{>b}\quad\text{si } a<b,\ a,b\in J .
\]

\struck{\ill{}} Dans le cas où $I=\{\omega_0,\omega_1\}$ \uncertain{coïncide} \ill{}, on a $J=\emptyset$, qu'il ne faut nullement exclure ! Les strates de $\mathcal{M}$ se réduisent au seul élément $[J]$. On trouve le magasin réduit à une seule strate, mais visualisé comme figure un \textit{\uncertain{intervalle ouvert indivisible}} : l'intervalle ouvert indivisible est représenté par l'ensemble ordonné vide !

Si on veut une formulation \uncertain{commune}, il faut dire qu'on se donne
\[
  (I,\Omega_0,\Omega_1)
\]
où
\[
  I,\ \text{un ordonné}; \qquad \Omega_0,\Omega_1\subset I
\]
et les conditions
\begin{enumerate}
  \item[a)] \struck{$\forall$} $x\in\Omega_0 \Rightarrow x$ est un plus petit élément de $I$
  \item[a')] $y\in\Omega_1 \Rightarrow y$ \textit{id.} plus grand \textit{id.}
  \item[b)] $x\in\Omega_0,\ y\in\Omega_1 \Rightarrow x\neq y$.
\end{enumerate}
\note{à la ligne a'), des traits de répétition (« idem ») encadrent « plus grand ».}
\marginal{donc $\operatorname{card}\Omega_0,\ \operatorname{card}\Omega_1\leq 1$, et il y a quatre cas à distinguer}

\page{26}
\note{p. 23 de l'auteur.}
Tant qu'il ne s'agit que de décrire un magasin, et y faire une théorie des polygones, des parties constructibles et des supports, \add{et \uncertain{constructibles}}, qui \uncertain{dominent}, il n'y a pas à faire d'autres hypothèses sur $J$.

Quand on regarde les polygones subordonnés à un diagramme, on \struck{peut prendre \ill{}} a intérêt (pour l'énoncé des résultats principaux) à prendre un diagramme qui contienne
\[
  \partial J \overset{\mathrm{df}}{=} \Omega_0\cup\Omega_1 .
\]

Il y a lieu de regarder la structure des \add{\ill{}} magasins constructibles \add{(\uncertain{ou mieux des supports})} induits sur des parties constructibles \add{$\Phi$} elles-mêmes, i.e.\ des polygones constructibles.

Voici le résultat :

a) $\mathcal{M}(\Phi)$ est la somme directe (à définir !) des $\mathcal{M}(\Phi_i)$, où les $\Phi_i$ sont les composantes \add{\ill{}} connexes de $\Phi$. Et \uncertain{un} \ill{} de la forme $\{a\}$, $a\in I$
\marginal{$\Phi=\{a\}$ cas trivial, $\mathcal{M}=$ \uncertain{strate} \ill{} ($a$ \ill{})}
\struck{$\mathcal{M}(\Phi$} Si $\Phi$ \add{\uncertain{connexe}}, ($\mathcal{M}(\Phi)$ est de la forme
\[
  \mathcal{M}(J,\Omega_0,\Omega_1), \quad\text{où } J,\Omega_0,\Omega_1
\]
sont définis ainsi, en termes de $\Phi$ :
\[
  (*)\quad J=\Sigma_\Phi\cap\mathcal{M}_0=\operatorname{Omb}^{\circ}_{I}(\Phi)
\]
\[
  =\Bigl\{s\in I \Bigm| \text{a) } a\leq s\leq b,\ \ \text{b) } s\cup\delta\Phi\in\operatorname{Diag}_I,\ \ \text{c) } s\in\{a,b\}\Rightarrow s\in\Phi^{\circ}\Bigr\}
\]
\marginal{$a=\operatorname{or}\delta\Phi$, $b=\operatorname{ex}\delta\Phi$}
\textit{Attention} cette description \ill{} si $\Phi$ \ill{}.
\note{bas de page serré, en partie au-dessous de la dernière ligne réglée ; lecture fragmentaire. La condition b) est lue « $s\cup\delta\Phi\in\operatorname{Diag}_I$ » sous réserve.}

\page{27}
\note{p. 24 de l'auteur.}
Donc, si à chaque strate $X$ \add{de $\mathcal{M}$} on associe son support intérieur $|X|^{\circ}\subset I$, \uncertain{que} la relation d'ordre induit \uncertain{sur} $I$ (NB $|X|^{\circ}$ peut être vide \add{si $X\in\mathcal{M}_1$} !), alors $J$ est la somme \add{lexicographique} \struck{ordonnée des} des
\[
  |X|^{\circ}\quad (X\in\Phi),
\]
\uncertain{en suivant} dans l'ordre $\trianglelefteq$ de $\mathcal{M}$.
\marginal{cette description \uncertain{lexicographique} \ill{}}
D'autre part, \struck{$\Omega_0$ est vide, ou} \struck{\ill{}}
\[
  \Omega_0=\begin{cases}
    \emptyset & \text{si la première strate de $\Phi$ n'a pas d'origine,}\\
              & \text{i.e. est un $I_{<a}$, ou si elle a une origine $a\notin\Phi$}\\
    \{a=\operatorname{or}(\Phi)\} & \text{sinon}
  \end{cases}
\]
\[
  \text{\struck{$\Omega_1=\operatorname{ex}(\Phi)\cap$}}
\]
$\Omega_1$ défini dualement${}^{*}$.
\note{dans l'accolade, une première ligne « $\operatorname{or}(\Phi)\cap\Phi$ » est biffée.}

\textit{NB} Il faut définir l'origine d'un polygone, quand elle existe (i.e.\ quand sa première strate \struck{\ill{}} n'est pas un $I_{<a}$), et dualement, l'extrémité.

\note{un trait horizontal sépare ici deux alinéas.}

Cas $\mathcal{M}$ \textit{divisible}. Dans le cas $(I,\omega_0,\omega_1)$ cela signifie que pour \struck{$a<b$} $a<b$ ($a,b\in I$) $\exists\, c$ avec $a<c<b$. Dans le cas de $(I,\omega_0)$ resp.\ $(I,\omega_1)$, il faut ajouter que $\forall a\in I$, existe $b\in I$ t.q.\ $b<a$ resp.\ $b>a$. Enfin, dans le cas de l'ensemble $I$ \uncertain{bi-ouvert}, il \add{faudrait} \struck{\ill{}} ajouter encore : $I\neq\emptyset$ (\textit{NB} pour dire que l'\textit{ensemble ordonné} $I$ est divisible : donc $\emptyset$ ne serait pas considéré comme un ordonné divisible, bien que $\mathcal{M}(\emptyset)=\{\emptyset\}$ \uncertain{soit} divisible, \uncertain{ce qui gêne} pourtant, à la \uncertain{suite} !)

Si $I$ est divisible, alors \uncertain{l'ens.} $\mathcal{M}_c\subset\mathcal{M}$ des strates « compactes » (i.e.\ de la forme $a$ et $I_{a,b}$) est très dense.

\page{28}
\note{p. 25 de l'auteur.}
\textit{Digression 1} \quad Somme directe de magasins.

Si $(\mathcal{M}_i)_{i\in I}$ est une famille de magasins, on définit le magasin somme $\mathcal{M}$ ainsi :
\begin{enumerate}
  \item[a)] $\mathcal{M}=\coprod_i\mathcal{M}_i$.
  \item[b)] $\trianglelefteq$ sur $\mathcal{M}$ est la relation d'ordre somme
  \item[c)] $\mathring{\ll}$ sur $\mathcal{M}$ est la relation d'ordre somme \ill{}
  \item[d)] si $x,y\in\mathcal{M}$, on a
  \[
    x\mathbin{|0|}y \iff \begin{cases}
      x\in\mathcal{M}_i,\ y\in\mathcal{M}_j,\ i\neq j\\
      \text{\textit{ou}}\\
      x,y\in\mathcal{M}_i,\ x\mathbin{|0|}y \text{ au sens de } \mathcal{M}_i .
    \end{cases}
  \]
\end{enumerate}
Vérification des axiomes tautologique.

Si pour chaque $i\in I$, on a choisi $\mathcal{F}_i\in\mathfrak{P}(\mathcal{M}_i)$, atelier associé à $\mathcal{M}_i$, on définit
\[
  \mathcal{F}\subset\mathfrak{P}(\mathcal{M}),\qquad \mathcal{F}=\{A\subset\mathcal{M}\mid \forall i\in\mathcal{M}_i,\ A\cap\mathcal{M}_i\in\mathcal{F}_i\}.
\]
On voit que c'est un atelier associé à $\mathcal{M}$, appelé \textit{atelier somme} des ateliers $(\mathcal{M}_i,\trianglelefteq,\mathring{\ll},|0|,\mathcal{F}_i)$. On notera
\[
  \mathcal{F}=\prod_{i\in I}\mathcal{F}_i .
\]
\uncertain{Dans} \ill{} \ill{} les ordres $\trianglelefteq$, $\subseteq$, $\ll$.
\note{« $\forall i\in\mathcal{M}_i$ » est sur la page ; on attendrait « $\forall i\in I$ ».}

\page{29}
\note{p. 26 de l'auteur.}
\marginal{\ill{} $|0|$ \ill{} (note en oblique dans la marge gauche, en haut)}
On fera attention que cette notion de somme \ill{} \uncertain{n'est pas} une notion de somme \uncertain{au sens} catégorique (pour la notion de \uncertain{(co)morphismes} de magasins \uncertain{introduite} jusqu'à présent). Ce ne serait \uncertain{pas vrai}, si l'on n'avait pas \uncertain{la condition} géométrique de \uncertain{raccordement} des flèches — \uncertain{intersection} ; la somme de magasins \uncertain{correspond} bien à une somme d'espaces. \uncertain{Cela tient} à un produit, \uncertain{qu'on pourrait} \ill{} la \uncertain{condition} principe de \uncertain{raccordement} des flèches).

Bien sûr \ill{} \uncertain{pour} un magasin $\mathcal{M}$, on \uncertain{peut le décomposer} : définir ses composantes connexes $\mathcal{M}_i$, comme les composantes connexes pour la relation $\ll$. On trouve une décomposition en somme, \uncertain{à condition de supposer} l'axiome suivant (qui sera \uncertain{bien évidemment} \uncertain{vérifié} dans le \uncertain{présent}) :

\textit{\uncertain{Mag connexe}} \quad Si $X,Y\in\mathcal{M}$ \uncertain{appartiennent} à des composantes connexes distinctes de $\mathcal{M}$, on a $X\mathbin{|0|}Y$.

\page{30}
\note{p. 27 de l'auteur.}
\textit{Digression 2} \quad Notion de sous-magasin.

Soit $\mathcal{M}$, \add{$\trianglelefteq,\mathring{\ll},|0|$,} un magasin, $\mathcal{M}'$ une partie, munie des relations induites ; quand peut-on dire, \uncertain{simplement}, que c'est un magasin.

\textit{Mag 1} trivial.

\textit{Mag 2} \uncertain{aussi}. Pour b, il faut supposer : \add{\textit{Sous-mag}} si $X\mathbin{\mathring{\ll}}Y$ dans $\mathcal{M}'$, et $X'\trianglelefteq X$ dans $\mathcal{M}'$, alors \add{\ill{}} l'\uncertain{existence} $Y'\in\widetilde{Y}$ tel que $X'\mathbin{\mathring{\ll}}Y'$ \uncertain{est} \add{\uncertain{automatique}} \uncertain{dans} $\mathcal{M}'$ \uncertain{(vérification automatique si $\mathcal{M}'$ fermé par $\trianglelefteq$)}.

\textit{Mag 3} trivial, \textit{Mag 4} trivial.

\marginal{NB Il faudrait imposer une condition plus forte, \ill{} \textit{Sous-mag} : si $X\mathbin{\mathring{\ll}}Y'\trianglelefteq Y$, et $X,Y\in\mathcal{M}'$, alors $Y'\in\mathcal{M}'$, laquelle \uncertain{équivaut} à \textit{Sous-mag} quand $\mathcal{M}'$ \ill{} \ill{} \struck{alors} $X,Y\in X$ \ill{} les exemples plus bas \uncertain{considérés} \uncertain{conservent} Mag 3.}
\note{note marginale écrite en oblique dans la marge gauche, de « Mag 1 » à « Ex 1 » ; lecture partielle.}

\textit{Ex 1} Soit $\Phi$ \add{$\subset\mathcal{M}$} un polygone de $\mathcal{M}$, i.e.\ une partie telle que $\overline{\Phi}^{(\trianglelefteq)}$ soit une figure, i.e.\ \uncertain{encore}, \struck{$\Phi$ fermé} $\forall X,Y\in\Phi$, on a $X\not\geq Y$. Considérons $\operatorname{Omb}^{\circ}(\Phi)$ (\uncertain{attention}, c'est fermé par $\mathring{\ll}$, mais pas en général par $\trianglelefteq$). \struck{Pour que} \add{Alors} $\mathcal{M}'=\operatorname{Omb}^{\circ}(\Phi)$ est un sous-magasin, i.e.\ satisfait Sous-mag \struck{\ill{}}.

\struck{soit un sous-magasin, i.e.\ \ill{} \ill{}. Condition Sous-mag ci-dessus, il faut et il suffit qu'on l'ait. Sous-mag $\Phi$ : $\forall Z\in\Phi$ et $X\mathbin{\mathring{\ll}}Z$, soit $X'\in\widetilde{X}$, et $Z'$ son image dans $\widetilde{Z}$ (donc $Z'\in\overline{\Phi}$), on veut que $X'\in\operatorname{Omb}^{\circ}(\Phi)\Rightarrow Z'\in\operatorname{Omb}^{\circ}(\Phi)$ i.e.\ en fait $Z'\in\Phi$ (puisque $Z'\in\overline{\Phi}$)} (\uncertain{même} s'il n'est pas fermé par $\trianglelefteq$, quand $\Phi$ n'est pas une figure…)
\note{tout le dernier alinéa, de « soit un sous-magasin » à « $Z'\in\overline{\Phi}$) », est barré de longs traits obliques et encadré à gauche ; la parenthèse finale reste hors de la rature.}

\page{31}
\note{p. 28 de l'auteur.}
En effet, soient $X,Y,X'\in\operatorname{Omb}^{\circ}(\Phi)$, avec
\[
  \begin{array}{ccc}
    X & \mathring{\ll} & Y\\
    \triangledown\!\!/ & & \\
    X' & &
  \end{array}
\]
alors $X'\ll Y\ll\overline{\Phi}$ donc $X'\ll\overline{\Phi}$, donc $\exists\,Z'_0\in\overline{\Phi}$ avec $X'\mathbin{\mathring{\ll}}Z'$. Comme $X'\in\operatorname{Omb}^{\circ}(\Phi)$, i.e.\ $\exists\,Z'_1\in\Phi$ tel que $X'\mathbin{\mathring{\ll}}Z'_1$, on a $Z'_1=Z'_0$ i.e.\ $Z'_0\in\Phi$. \uncertain{Ceci posé}, il \uncertain{revient} \uncertain{au} \uncertain{lemme} $Y\in\operatorname{Omb}^{\circ}(\Phi)$ \uncertain{avec} $Z\in\Phi$, $Y\mathbin{\mathring{\ll}}Z$, d'où
\[
  \begin{array}{ccccc}
    X & \mathring{\ll} & Y & \mathring{\ll} & Z\in\Phi\\
    \triangledown\!\!/ & & \triangledown\!\!/ & & \triangledown\!\!/\\
    X' & \mathring{\ll} & Y' & \mathring{\ll} & Z'\\
     & & \overset{?}{\in} & & \\
     & & \operatorname{Omb}^{\circ}(\Phi) & &
  \end{array}
  \quad ;\ \text{et comme } Z'\in\overline{\Phi},
\]
donc $Z'=Z'_0\in\Phi$, donc $Y'\in\operatorname{Omb}^{\circ}(Z')\subset\operatorname{Omb}^{\circ}\Phi$, cqfd.
\note{le signe $\triangledown\!\!/$ rend le $\trianglelefteq$ couché qu'il écrit verticalement entre deux lignes ($X'\trianglelefteq X$, etc.).}

\textit{Ex 2} Soit $\Phi$ comme précédemment, et soit
\[
  \mathcal{M}'=\operatorname{supp}^{\circ}\Phi \quad(\text{un support constructible}).
\]
Est-ce un sous-magasin ? C'est pratiquement \uncertain{équivalent} à la même question, \uncertain{pour} $\mathcal{M}'=\operatorname{cosupp}^{\circ}\Phi$. \struck{\ill{} \ill{}, \ill{}} C'est faux en général, comme on voit sur la figure suivante.

\drawing{un triangle de sommets $a$ (en haut), $b$, $c$ ; la moitié gauche hachurée, un segment de $a$ au point $d$, marqué d'un gros point, sur le côté $[b,c]$.}
\[
  \begin{aligned}
    &Y=\text{2-simplexe }(a,b,c)\\
    &X=\text{2-simplexe }(a,b,d)\\
    &X'=\text{1-simplexe }(b,d)\\
    &\text{\struck{$\Phi=\ldots$}}\\
    &Y'=\text{1-simplexe }(b,c)\\
    &\Phi=\{d\}.
  \end{aligned}
  \qquad
  \text{On a}\quad \Phi\mathbin{|0|}X,\ \ \Phi\mathbin{|0|}Y,\ \text{et}\ \Phi\mathbin{|0|}X',
\]
\note{la ligne biffée sous $X'$ est illisible ; dans la ligne qui suit $\Phi=\{d\}$, le deuxième $|0|$ est lu sous réserve. L'argument continue à la page suivante.}

\page{32}
\note{p. 29 de l'auteur.}
mais on n'a pas $\Phi\mathbin{|0|}Y'$. Donc, attention à la définition du magasin visiblement \add{$F$} d'une figure (ici $\{d\}$, figure ponctuelle !).
\note{les trois lignes ci-dessus sont marquées d'un triple trait vertical dans la marge gauche.}

Il faut, \uncertain{si on} \uncertain{veut} se borner aux \uncertain{multi-}strates $X$ avec $X\mathbin{||}F$, \uncertain{qui est stable} par $\trianglelefteq$ \uncertain{donc un sous-magasin}, \textit{se borner aux $X$ telles que $X\not\geq F$}, ou $X\not\geq\Phi$ dans le cas d'un polygone. (NB dans le cas de \uncertain{la dimension} 1, \uncertain{exemple} \ill{}, on \uncertain{avait} $X\mathbin{|0|}\Phi\Rightarrow X\not\geq\Phi$, ou $X\mathbin{|0|}Y\Rightarrow X\not\geq Y$ …). \uncertain{Je dis} \add{plus \uncertain{précisément}} \uncertain{que si} \add{\uncertain{grand}}
\[
  \mathcal{M}'=\{X\in\mathcal{M}\mid X\mathbin{|0|}\Phi,\ X\not\geq\Phi\} \quad ?
\]
\marginal{magasin \uncertain{visible} hors de $\Phi$.}
alors $\mathcal{M}'$ est un sous-magasin, i.e.
\[
  \begin{array}{ccc}
    X & \mathring{\ll} & Y\\
    \triangledown\!\!/ & & \triangledown\!\!/\\
    X' & \mathring{\ll} & Y'
  \end{array}
  \qquad \text{alors}\quad X,Y,X'\in\mathcal{M}\Rightarrow Y'\in\mathcal{M}.
\]
Comme $Y\not\geq\Phi$ et $Y'\trianglelefteq Y$, on a $Y'\not\geq\Phi$. Donc pour avoir $Y'\mathbin{|0|}\Phi$, il suffit de voir $Y'\notin\Phi$. Mais si on avait $Y'\in\Phi$, on aurait $X'\mathbin{|0|}Y'$,
et comme $X'\mathbin{\mathring{\ll}}Y'$, $X'\mathbin{\|}X'$, absurde.
\note{au lieu de $X'\mathbin{\|}X'$ la page porte un signe surchargé entre deux barres verticales ; la lecture est douteuse. « $\mathcal{M}$ » dans « $X,Y,X'\in\mathcal{M}\Rightarrow Y'\in\mathcal{M}$ » est sur la page, où l'on attendrait $\mathcal{M}'$.}

\page{33}
\note{p. 30 de l'auteur.}
\marginal{\struck{8 juillet}}
\textit{Digression 3} \quad Magasins avec donnée de strates \add{dites} \textit{compactes} (\textit{Magasins compactologiques}).
\textit{Notion de figure localement finie}.

On se donne une partie $\mathcal{M}_c$ de $\mathcal{M}$, \add{\ill{}} \struck{\ill{}} (\uncertain{structure supplémentaire}) formée des strates dites \struck{« admissibles »} « \textit{compactes} ».

\struck{Ax} L'axiome préalable est le suivant :

\textit{Mag cpct 1} \quad \struck{a) Si $X$ cpct, \ill{} alors $\widetilde{X}$ fini.}
\begin{enumerate}
  \item[a)] $X\in\mathcal{M}_c\Longrightarrow\widetilde{X}$ fini
  \item[b)] $X\in\mathcal{M}_c,\ Y\ll X\Longrightarrow Y\in\mathcal{M}_c$ — en d'autres termes : $\mathcal{M}_c$ est fermé dans $(\mathcal{M},\ll)$, \uncertain{ou} \uncertain{ce qui revient au même}, dans $(\mathcal{M},\mathring{\ll})$ et $(\mathcal{M},\trianglelefteq)$.
\end{enumerate}

Par la suite, on sera amené sûrement à introduire un axiome \textit{Mag cpct 2}, \uncertain{de type} Borel-Lebesgue ; et \uncertain{aussi} que $\mathcal{M}_c$ est « riche » dans $\mathcal{M}$ (\uncertain{grâce aux subdivisions}).

Quand on se donne de plus un \add{$\mathcal{F}$} atelier, alors \uncertain{c'est} un magasin, \uncertain{on exige} :

\struck{\textit{Mag cpct (at)}}

\textit{Mag at cpct} \quad Si $F\ll X$, avec $F\in\mathcal{F}$, $X\in\mathcal{M}_c$, alors $\widetilde{F}$ fini. Mieux : si $X\in\mathcal{M}_c$, alors $\forall F\in\mathcal{F}$, l'ensemble des \struck{$X$} \add{$Y$} $\in\widetilde{F}$ telles que \struck{$\ldots|0|X$} $Y\mathbin{||}X$, i.e.\ $\widetilde{Y}\mathbin{|0|}\widetilde{X}$, a un complémentaire fini dans $\widetilde{F}$.

\struck{On notera \add{ceci : une compactologie donnée, \ill{}}. Pour $\mathcal{M}_c$ donné, il existe un \uncertain{$\mathcal{F}$} \uncertain{finissant} : $\mathcal{M}$, le plus grand possible, parmi ceux satisfaisant à la condition \uncertain{envisagée}. Les \uncertain{\ill{}} figures qui}
\note{les cinq dernières lignes de la page sont barrées de traits obliques ; la phrase s'interrompt au bas de la page.}

\marginal{Un magasin compactologique \ill{} est dit « à strates compactes » si $\mathcal{M}_c=\mathcal{M}$. Il est dit \textit{compact} si \ill{} $\exists F\in\mathcal{F}(\mathcal{M})$, avec $F$ \textit{fini}, tel que $\operatorname{supp}^{\circ}F=\mathcal{M}$, i.e.\ $\nexists X\in\mathcal{M}$, $F\mathbin{|0|}X$. … \textit{Figures} \struck{compactes} : $\widetilde{F}$ fini, et $\subset\mathcal{M}_c$. \textit{Polygones} \uncertain{compacts} : \struck{\ill{}} les figures compactes \ill{} \ill{}}
\note{note marginale en oblique, très serrée, dans la marge gauche ; lecture partielle.}

\page{34}
\note{p. 31 de l'auteur.}
Pour $\mathcal{M}_c$ donné, et une partie $\Phi\subset\mathcal{M}$, on dira que $\Phi$ est \textit{localement finie} si \struck{$\forall X\in\mathcal{M}_c$,}
\[
  \forall X\in\mathcal{M}_c,\ \text{\struck{posant}}\ \text{\add{l'ens.}}\ \Phi_X=\{Y\in\Phi\mid Y\mathbin{\overline{\|}}X\}\ \text{\add{est fini}}
\]
\[
  \text{\struck{, on a $\Phi\smallsetminus\Phi_X$ fini.}}
\]
Les parties loc.\ finies de $\mathcal{M}$ forment un \uncertain{anti-filtre}. L'ensemble \struck{des} $\mathcal{F}_{\mathrm{\ell f}}(\mathcal{M})$ des figures loc.\ finies de $\mathcal{M}$ forme un \add{$\mathcal{M}$-}atelier, à condition \uncertain{de} \uncertain{partir} des \uncertain{supports} \struck{qui} \uncertain{ici}. \uncertain{On a}
\[
  \mathcal{M}\subset\mathcal{F}_{\mathrm{\ell f}}(\mathcal{M})
\]
\note{la page porte bien $\mathcal{M}\subset\mathcal{F}_{\mathrm{\ell f}}(\mathcal{M})$ ; il faut sans doute entendre que chaque strate, vue comme figure, est localement finie.}

\textit{Mag cpct 1} c) \struck{\ill{}} $X\in\mathcal{M}$, $\widetilde{X}$ est loc.\ fini pour $\mathcal{M}_c$ i.e.\ $\forall Y\in\mathcal{M}_c$, l'ensemble des
\[
  \{X'\in\widetilde{X}\mid X'\mathbin{\overline{\|}}Y\}
\]
est fini.

Si $\Phi$ est un polygone loc.\ fini, alors $\overline{\Phi}^{(\trianglelefteq)}$ est encore loc.\ fini.

L'axiome \textit{Mag at cpct} signifie que $\mathcal{F}\subset\mathcal{F}_{\mathrm{loc\,fin}}$. Ainsi $\mathcal{F}_{\mathrm{loc\,fin}}$ est le plus grand atelier \uncertain{définissant} \uncertain{le} $\mathcal{M}$, qui soit compatible à la compactologie.

Inversement, donnons-nous
\[
  (\mathcal{M},\trianglelefteq,\mathring{\ll},|0|,\mathcal{F}).
\]
On va construire une \struck{sous} compactologie la plus grande possible, pour laquelle $\mathcal{M}_c$, $\mathcal{F}$ soient compatibles. On posera
\[
  \mathcal{M}_c=\{X\in\mathcal{M}\mid \forall F\in\mathcal{F},\ \text{l'ens. } \{Y\in\widetilde{F}\mid X\mathbin{\overline{\|}}Y\}\ \text{est fini}\}
\]
\note{la relation $\overline{\|}$ (deux barres verticales surmontées d'un trait) est transcrite telle qu'elle est dessinée ; à la page précédente, la même relation est écrite sans trait, $\|$, et glosée « i.e. $\widetilde{Y}\mathbin{|0|}\widetilde{X}$ ».}

\page{35}
\note{p. 32 de l'auteur.}
Il est clair que pour cet $\mathcal{M}_c$, les conditions \textit{Mag cpct 1} a) b) \struck{c)} \add{et c)} sont satisfaites (pour c), \uncertain{on} prend $F=\widetilde{X}$).

\struck{Ab} Si $\mathcal{F}$ est formé de figures finies, \uncertain{tout} \uncertain{fait} $\mathcal{M}_c$ est aussi la plus grande des compactologies sur $\mathcal{M}$. Mais il ne semble pas raisonnable de \textit{définir} en général $\mathcal{M}_c$ en termes des seuls magasins $(\mathcal{M},\trianglelefteq,\mathring{\ll},|0|)$, par les conditions précédentes, comme \struck{les}
\[
  \mathcal{M}_c=\{X\in\mathcal{M}\mid \forall Y\in\mathcal{M},\ \text{l'ens. } \{Y'\in\widetilde{Y}\mid Y'\mathbin{\overline{\|}}X\}\ \text{fini}\},
\]
lequel donnerait $\mathcal{M}_c=\mathcal{M}$ quand \uncertain{toutes} les multistrates \struck{de} $\widetilde{Y}$ de $\mathcal{M}$ sont finies. P.\,ex.\ dans le cas d'un ensemble ordonné $I$, muni d'une seule extrémité \uncertain{ou} \uncertain{sans} extrémité (\textit{intervalles ouverts} qui n'est pas un segment), il est naturel de considérer comme compactes les seules strates de la forme $a$ ($a\in I$) $I_{a,b}$ ($a,b\in I$).

Il semble bien que ce soit la donnée supplémentaire de la compactologie, qui soit une structure \uncertain{souple} plus intéressante que celle d'un atelier \uncertain{tout} particulier.

\page{36}
\note{p. 33 de l'auteur.}
Elle permet, en plus des \struck{ateliers des} figures compactes (et des figures finies dont on disposait déjà avant), d'introduire les figures \textit{localement finies}. Il ne semble pas qu'en pratique, il y ait lieu de travailler avec d'autres ateliers attenants que celui des figures finies, celui des figures \textit{localement finies}, ou (à la rigueur) celui de \textit{toutes} les $\mathcal{M}$-figures.

\textit{NB} Les figures \textit{compactes} forment un atelier attenant à $\mathcal{M}_c$, qui est un \uncertain{sous-atelier} de $\mathcal{M}$, mais non à $\mathcal{M}$, \uncertain{sauf} si $\mathcal{M}_c=\mathcal{M}$.

Ainsi, il \uncertain{paraît} \uncertain{clair} que la bonne \add{(tautologique)} \uncertain{approche} des \uncertain{considérations} de la topologie « \uncertain{supportique} » \ill{} \ill{} \uncertain{que je propose}, est plutôt \struck{la notion} que la notion de \textit{magasin compactologique}
\[
  (\mathcal{M},\trianglelefteq,\mathring{\ll},|0|,\mathcal{M}_c).
\]
C'est une donnée plus simple que celle d'un atelier attenant à $\mathcal{M}$, \uncertain{ou} $\mathcal{M}_c\in\mathfrak{P}(\mathcal{M})$, et $\mathcal{F}\in\mathfrak{P}(\mathfrak{P}(\mathcal{M}))$.
\note{la fin de la phrase est lue « où $\mathcal{M}_c\in\mathfrak{P}(\mathcal{M})$, et $\mathcal{F}\in\mathfrak{P}(\mathfrak{P}(\mathcal{M}))$ » : la comparaison oppose une partie de $\mathcal{M}$ à une partie de $\mathfrak{P}(\mathcal{M})$. Les mots entre « topologie » et « est plutôt » sont pour la plupart douteux.}

\page{37}
\note{p. 34 de l'auteur.}
\textit{Digression 4} \quad \add{Terminologie : atelier, magasin, strates, multistrates.} (Ainsi, de plus en plus, l'accent principal dans le travail de fondements, se porte sur $\mathcal{M}$, plutôt que sur $\mathcal{F}$. Du coup, cela suggère aussi un changement de terminologie : le \struck{terme} \add{nom} initial principal, d'\textit{atelier}, \textit{devrait être donné à $\mathcal{M}$}. Cet atelier travaille avec les objets qui figurent dans $\mathcal{M}$, les \add{strates ou} multistrates, pour constituer avec eux des \textit{figures}. L'ensemble de ces figures peut-être considéré comme constituant un \textit{magasin attenant} : l'atelier (plutôt que l'inverse) — plusieurs magasins pouvant être attenants au même atelier, comme on \uncertain{vient} de le \uncertain{voir} \uncertain{plus haut}. Dans le cas d'un \uncertain{atelier} compactologique, on pourrait \uncertain{appeler} peut-être $\mathcal{M}_c$, formé des (multi)strates compactes, le \textit{cœur} de l'atelier.
\note{il propose ici d'échanger les noms : « atelier » pour $\mathcal{M}$, « magasin attenant » pour l'ensemble $\mathcal{F}$ des figures. La terminologie des pages précédentes (« magasin » pour $\mathcal{M}$, « atelier » pour $\mathcal{F}$) est transcrite telle qu'elle est écrite. La parenthèse ouverte après « Digression 4 » n'est pas refermée sur la page.}

\page{38}
\note{p. 35 de l'auteur.}
Il \uncertain{vient} \uncertain{la} perplexité sur la terminologie parallèle strates — multistrates. Finalement, je vais appeler les $X\in\mathcal{M}$ les \textit{strates} de l'atelier $\mathcal{M}$. Chaque strate détermine l'ensemble $\widetilde{X}$ de ses \textit{sous-strates}, ou \textit{strates incidentes}, C'est $\widetilde{X}$ qui serait la \textit{multistrate associée} à la strate $X$, ou détermination, définie par $X$, \struck{ou la}

Ainsi, pour une figure $F$, on distingue entre les \textit{strates de} $F$, qui sont les éléments de $\widetilde{F}$, ou \textit{strates de $\mathcal{M}$ incidentes} à $F$, et les \textit{multistrates de $F$}, qui sont les $\widetilde{X}$ pour $X\in\widetilde{F}$. On pourrait désigner l'ens.\ de ces multistrates par $\widetilde{\widetilde{F}}$.

Quand on décrit l'atelier des figures \uncertain{définissant} \struck{les} un ensemble \add{$\mathcal{L}$}
\[
  \operatorname{Figil}(\mathcal{L})=\mathcal{M}
\]
il est entendu que \uncertain{ces} objets apparaissent \add{\uncertain{donc}} non comme des parties, mais comme des ensembles de parties de $\mathcal{L}$, \struck{soit} $\widetilde{X}$. Cela \uncertain{n'empêche} qu'une telle ensemble de parties de $\mathcal{L}$, considéré comme objet $\in\mathcal{M}$, \uncertain{sera} appelée « strate », et non « multistrate ». Ainsi
\note{« $\operatorname{Figil}$ » est lu sous réserve. La phrase « Ainsi » continue à la page suivante.}

\page{39}
\note{p. 36 de l'auteur.}
alors qu'un objet $X$ de $\mathcal{M}$ est un \struck{\ill{}} élément (de nature particulière) de $\mathfrak{P}(\mathfrak{P}(\mathcal{L}))$, une multistrate $\widetilde{X}$ de $\mathcal{M}$ doit être considérée, en bonne logique, comme un élément de $\mathfrak{P}(\mathfrak{P}(\mathfrak{P}(\mathcal{L})))$ — un ensemble d'ensembles de parties de $\mathcal{L}$.

De plus en plus, $\mathcal{M}$ apparaissant comme l'ensemble de base privilégié, les \textit{figures} $F$ apparaissent comme des parties de $\mathcal{M}$. Dans ce point de vue, il n'y a donc pas lieu d'introduire un $\widetilde{F}$, qui ici n'est autre que $F$. Donc il n'y a pas lieu de parler de « déploiement » d'une figure. C'est à l'avenir qu'on verra s'il y a lieu, par la suite, d'autonomiser $\mathcal{F}$ p.r.\ à $\mathcal{M}$, par l'introduction d'une structure mixte sur $(\mathcal{M},\mathcal{F})$, avec une relation d'incidence $\trianglelefteq$ entre $\mathcal{M}$ et $\mathcal{F}$ — auquel cas la notation $\widetilde{F}$ reprend un sens et ses droits.

\page{40}
\note{p. 37 de l'auteur.}
Dès maintenant, je dirai « atelier » là où je disais « magasin », et inversement. Les axiomes Mag 1 à Mag 4 deviennent désormais At 1 à At 4, etc.

Je reviens à la notion de supports (« contours ») associés à un atelier — que je \uncertain{ne} vais plus \uncertain{noter} $\mathcal{M}$, il faut revoir les \uncertain{conventions} de ses notations — mais $\mathcal{A}$, donc
\[
  (\mathcal{A},\trianglelefteq,\mathring{\ll},|0|).
\]
La donnée de $|0|$, relation symétrique antiréflexive, nous permet de définir un ensemble de parties particulières de $\mathcal{A}$, appelées « supports », parties fermées par $\mathring{\ll}$ (mais non par $\trianglelefteq$ en général), \struck{la donnée}
\[
  \Sigma_{\mathcal{A}}\subset\mathfrak{P}(\mathcal{A}),
\]
stable par intersections quelconques, qui
\note{la phrase continue au-delà de ce lot.}

\end{document}
